\section{Resultados del experimento y análisis}

\subsection{Comparación de FID bajo diferentes valores de $\eta$ y número de pasos}

El artículo sobre DDIM realizó un ablitivo experimental sistemático en el conjunto de datos CIFAR-10, evaluando el impacto de $\eta$ y el número de pasos de muestreo $S$ sobre el FID.

Las observaciones clave son las siguientes:

\begin{enumerate}
    \item \textbf{Todos los pasos ($S=1000$)}: Los valores de FID para todos los $\eta$ son bastante similares, aproximadamente 3-4. DDPM ($\eta=1$) es ligeramente mejor que DDIM determinista ($\eta=0$).
    \item \textbf{Número medio de pasos ($S=50\sim100$)}: DDIM determinista ($\eta=0$) supera significativamente a DDPM ($\eta=1$), con una diferencia en FID que puede alcanzar 5-10.
    \item \textbf{Pocos pasos ($S=10\sim20$)}: DDIM determinista aún mantiene una buena calidad (FID $\sim$10-15), mientras que el FID de DDPM se degrada drásticamente ($>$50).
\end{enumerate}

\begin{warningbox}{Compensación entre calidad y velocidad al reducir los pasos}
Aunque DDIM determinista es más robusto al reducir el número de pasos, la calidad disminuirá inevitablemente si los pasos bajan demasiado. Esto se debe a que un tamaño de paso temporal mayor provoca saltos más grandes en cada actualización, lo que reduce la precisión de la red de predicción del ruido en regiones alejadas de la distribución de entrenamiento. Los experimentos del artículo sobre DDIM demuestran que aproximadamente $S=50$ pasos en CIFAR-10 logran una calidad comparable a los 1000 pasos de DDPM.
\end{warningbox}

\subsection{Observaciones cualitativas de la calidad de las muestras}

El ponente también discutió algunas características cualitativas de las muestras generadas por DDIM:

\begin{itemize}
    \item Las imágenes generadas por DDIM determinista suelen ser ligeramente más ``nítidas'' (sharper) que las de DDPM, debido a la ausencia de inyección de ruido adicional.
    \item La aleatoriedad de DDPM puede aportar, en ocasiones, mayor diversidad y una sensación de ``naturalidad''.
    \item Cuando hay suficientes pasos ($\geq$50), la diferencia en la calidad visual entre ambos es mínima.
\end{itemize}

\subsection{Resumen de la sección}

Los resultados del experimento validan la ventaja de DDIM al reducir el número de pasos de muestreo: la degradación de la calidad con una muestra determinista es más lenta al disminuir los pasos. Aproximadamente 50 pasos de DDIM pueden igualar la calidad de 1000 pasos de DDPM, logrando una aceleración de aproximadamente 20 veces.
